% ============================================================
% Appendix Z.11 — FRFP Case Study (Autonomy / Mobility):
% Waymo Driverless Ride-Hailing
% ============================================================

\section{FRFP Case Study (Autonomy / Mobility): Waymo Driverless Ride-Hailing}
\label{app:frfp-waymo}

This appendix analyses Waymo’s driverless ride-hailing service as an FRFP
case study in long-horizon autonomy.
We focus on how its technology, safety processes, and deployment model map
onto FRFP’s explicit/tacit split, boundary design, safe horizons, and
institutional governance.

\subsection{Background}

Waymo, originally the Google self-driving car project (launched in 2009),
has evolved into a commercial provider of robotaxi services.
As of 2025, Waymo operates fully driverless ride-hailing (no human safety
driver in the car) in multiple U.S.\ cities, including Phoenix, San
Francisco, and parts of Los Angeles and other regions.\footnote{For an
overview of Waymo’s deployments and miles driven, see public documentation
and safety reports on the Waymo and Alphabet websites, as well as regulatory
filings and independent analyses.}

Key milestones (very briefly):

\begin{itemize}
  \item Early closed-course and supervised public-road testing (2010s);
  \item Limited public pilot in Chandler, Arizona, with safety drivers;
  \item Transition to fully driverless operations (no safety driver) in
        parts of the Phoenix metro area, expanding to 24/7 service;
  \item Gradual expansion to San Francisco and Los Angeles, with
        driverless service in designated operating domains;
  \item Public safety reports and disengagement statistics through
        California and Arizona regulators, and Waymo’s own safety cases.
\end{itemize}

Compared to many other high-profile autonomy efforts, Waymo has:

\begin{itemize}
  \item accumulated very large numbers of driverless miles without a
        correspondingly large number of serious crashes;
  \item maintained a comparatively cautious, constrained operational domain;
  \item invested heavily in explicit safety cases and institutional
        relationships with regulators and local authorities.
\end{itemize}

From an FRFP perspective, Waymo is a prominent example of an AI-intensive,
agentic system that has reached real-world, at-scale deployment in a domain
with stringent safety demands.

\subsection{Explicit and Tacit Spaces in Waymo’s System}

\subsubsection*{Explicit space \(E_{\mathrm{Waymo}}\)}

Waymo’s explicit domain consists of the artefacts processed by its
autonomous driving stack and operations:

\begin{itemize}
  \item \textbf{Perception data and models}:
        sensor inputs (LiDAR, radar, cameras), object detections, semantic
        segmentations, and tracking.
  \item \textbf{HD maps and localisation}:
        high-definition road maps, lane-level geometry, traffic control
        information, static infrastructure features.
  \item \textbf{Prediction and planning}:
        models of other agents’ behaviour; trajectory plans; control
        commands (steering, throttle, braking).
  \item \textbf{Operational data}:
        \begin{itemize}
          \item logs of each trip (telemetry, interventions, incidents);
          \item simulation scenarios and replay environments;
          \item safety metrics (e.g.\ collision rates, policy violations).
        \end{itemize}
  \item \textbf{Policies and procedures}:
        formal documentation of safety protocols, fallback behaviours,
        disengagement rules, and operational design domains (ODDs).
\end{itemize}

These are all explicit artefacts in \(E_{\mathrm{Waymo}}\), manipulated by
Waymo’s software systems and interpreted by engineers and regulators.

\subsubsection*{Tacit space \(T_{\mathrm{mobility}}\)}

Tacit space contains the epistemic and normative states of:

\begin{itemize}
  \item safety drivers (in earlier phases) and remote operations staff:
        their sense of when the system is behaving oddly or unsafely;
  \item Waymo’s safety engineers and leadership:
        judgement about acceptable risk, release readiness, and go/no-go
        decisions for new features and cities;
  \item regulators and city officials:
        beliefs about safety, public acceptability, and regulatory
        obligations;
  \item passengers and the public:
        evolving trust, expectations, and perceptions of risk vs.\ benefit.
\end{itemize}

Correctness in this domain is not just “did the car follow the rules of the
road?” but “is this system, in this city, under these conditions, safe
enough to operate driverlessly?”, which is a tacit institutional judgement.

\subsection{Boundary Design:
  From Onboard Autonomy to Human Oversight}

\subsubsection*{Onboard boundary: the vehicle as an explicit agent}

Within the vehicle, the autonomy stack operates as an explicit agent:

\begin{itemize}
  \item It perceives the environment, predicts others’ behaviour, and plans
        and executes a trajectory in real-time.
  \item In driverless service, there is no human in the loop for routine
        control; the system is authorised to actuate steering, acceleration,
        and braking directly.
  \item Safety fallbacks (e.g.\ minimal-risk conditions, safe stop) are
        encoded explicitly and triggered when the system detects
        uncertainties or failures.
\end{itemize}

In FRFP terms:

\begin{itemize}
  \item many low-level and mid-level decisions are fully automated within
        \(N \ltimes E_{\mathrm{Waymo}}\);
  \item correctness for these decisions is delegated to explicit semantics
        (control algorithms, safety monitors, redundancy mechanisms).
\end{itemize}

\subsubsection*{System-level boundary: fleet operations and governance}

Above the vehicle level, there are stronger boundaries:

\begin{itemize}
  \item \textbf{Operational Design Domain (ODD)}:
        driverless operation is restricted to predefined regions, weather
        conditions, and scenarios where the system has been validated.
  \item \textbf{Fleet operations centre}:
        remote staff can monitor vehicles, respond to unusual events, and
        provide high-level guidance or issue commands (e.g.\ instructing
        the vehicle to pull over and stop).
  \item \textbf{Safety release process}:
        major expansions (e.g.\ new cities, night driving, highway segments)
        go through internal review and, in some jurisdictions, explicit
        regulatory approval.
\end{itemize}

These structures place the key correctness boundary at:

\begin{itemize}
  \item decisions about where and when the system may operate driverlessly;
  \item how incidents and anomalies are handled at the fleet and governance
        level.
\end{itemize}

This matches FRFP’s principle that high-level correctness for long-horizon,
high-stakes systems must remain with human and institutional actors in
\(T_{\mathrm{mobility}}\), even if many sub-decisions are fully automated.

\subsection{Protocol-Level Properties in FRFP Terms}

\subsubsection*{Z.11.1 Explicit Semantics and Strong Evaluation}

Compared to many AI systems, autonomous driving has relatively strong,
structured explicit semantics:

\begin{itemize}
  \item Rules of the road, traffic regulations, and physical constraints
        define a clear set of admissible behaviours at the vehicle level.
  \item Safety metrics (e.g.\ crashes per million miles, policy-violation
        rates, near-miss statistics) can be defined and tracked at scale.
  \item Simulation and replay allow systematic stress-testing of the
        autonomy stack against rare or adversarial scenarios.
\end{itemize}

For Waymo:

\begin{itemize}
  \item large-scale data collection and years of testing have created rich
        explicit datasets; 
  \item incident and disengagement logs support post-hoc evaluation and
        refinement of explicit pipelines;
  \item safety reports and external analyses provide some transparency on
        performance and limitations.
\end{itemize}

In FRFP language, the mapping from sensor data and maps to control actions
has a well-defined semantics constrained by physics and law, and there is an
ongoing, explicit evaluation apparatus to test it.

\subsubsection*{Z.11.2 ODD as an Explicit Safe Horizon}

Waymo’s use of Operational Design Domains (ODDs) is a concrete instance of
FRFP’s safe-horizon concept:

\begin{itemize}
  \item The system is only allowed to operate driverlessly in regions,
        times, and conditions where it has been validated (e.g.\ specific
        parts of Phoenix and SF, typical weather, certain speed limits).
  \item Expansion of the ODD is gradual and gated:
        new neighbourhoods, nighttime operation, or more complex
        intersections are added only after extensive testing.
  \item When conditions fall outside the ODD (e.g.\ extreme weather,
        severe construction changes), the system may restrict service,
        request remote assistance, or execute safe stops.
\end{itemize}

From an FRFP perspective:

\begin{itemize}
  \item the ODD defines an explicit horizon in state-space:
        outside this region, the autonomy stack is not authorised to
        operate without additional tacit oversight;
  \item this enforces a structural limitation on how far explicit autonomy
        can run before requiring institutional re-grounding (via testing,
        recalibration, and regulatory review).
\end{itemize}

\subsubsection*{Z.11.3 Incident Handling and Feedback Loops}

Waymo’s operational model incorporates explicit feedback mechanisms:

\begin{itemize}
  \item Any incident or unusual behaviour (e.g.\ road blockage, confused
        negotiation with other vehicles, passenger complaints) is logged.
  \item Engineers can replay sensor data in simulation, examine why the
        behaviour occurred, and adjust perception, prediction, or planning
        modules.
  \item Policy changes (e.g.\ more conservative behaviour at certain
        unprotected left turns) can be rolled out fleet-wide.
\end{itemize}

In FRFP terms:

\begin{itemize}
  \item each incident creates a run segment in \(N \ltimes E_{\mathrm{Waymo}}\)
        that is then subjected to tacit evaluation by safety teams;
  \item this evaluation updates both explicit pipelines (software) and
        tacit states (risk understanding, operating procedures);
  \item correctness is thus enforced at the institutional boundary, with
        explicit artefacts feeding into a dynamic governance loop.
\end{itemize}

\subsection{Institutional and Societal Semantics}

\subsubsection*{Regulatory and city-level governance}

Waymo operates within a complex regulatory landscape:

\begin{itemize}
  \item State motor-vehicle regulators and public-utilities commissions
        grant permits for testing and driverless deployment.
  \item City authorities influence where vehicles may pick up and drop off,
        and how they interact with emergency services and road closures.
  \item Incidents are scrutinised by regulators, city officials, and the
        public, shaping tacit perceptions of safety.
\end{itemize}

FRFP views these institutions as operators on the population-level tacit
state \(T^{\mathrm{pop}}_{\mathrm{mobility}}\):

\begin{itemize}
  \item approvals, suspensions, and rule changes update beliefs about
        acceptable risk and control;
  \item public hearings and media coverage feed back into Waymo’s own safety
        posture and expansion plans;
  \item over time, a shared tacit understanding of what “safe enough”
        autonomy means in different cities is gradually constructed.
\end{itemize}

\subsubsection*{Public trust and behavioural adaptation}

Passengers and other road users form tacit beliefs about Waymo vehicles:

\begin{itemize}
  \item riders experience the cars as cautious but competent; many report
        high trust after a few trips, while others remain sceptical;
  \item human drivers learn to anticipate Waymo behaviour (e.g.\ its
        caution at narrow passes or merges);
  \item incidents---even rare ones---can cause sharp updates in public
        opinion, prompting calls for stricter oversight or pauses.
\end{itemize}

These dynamics matter for FRFP because they shape:

\begin{itemize}
  \item how often and where people choose to ride driverless vehicles;
  \item how regulators balance innovation vs.\ safety;
  \item how institutions design future autonomy policies and infrastructures.
\end{itemize}

\subsection{Counterfactual Pitfalls and FRFP Alignment}

Waymo’s trajectory can be contrasted with hypothetical or real failure modes:

\begin{itemize}
  \item \textbf{Overly broad ODD}.
        A less cautious system might claim nationwide operation in all
        conditions, leading to frequent boundary violations and unsafe
        behaviours.
        Waymo’s constrained ODDs are a partial FRFP-aligned safeguard.
  \item \textbf{Weak incident feedback}.
        Failing to systematically log and analyse incidents would erode
        the link between explicit artefacts and tacit safety judgements.
        Waymo invests heavily in these loops.
  \item \textbf{Purely metric-driven governance}.
        Over-reliance on disengagement statistics or proprietary safety
        metrics, without robust external scrutiny, would be an instance of
        explicit-only governance.
        Waymo’s external reporting and regulatory engagement, while not
        perfect, move toward a more FRFP-consistent institutional layer.
\end{itemize}

Nonetheless, FRFP would highlight ongoing open questions:

\begin{itemize}
  \item How well do current metrics capture tail risks and rare scenarios?
  \item Are city-level ODD boundaries and fleet behaviours transparent and
        contestable enough for the public and regulators?
  \item How should responsibility be allocated between Waymo, regulators,
        and riders in edge cases and accidents?
\end{itemize}

These questions sit squarely in tacit and institutional semantics, beyond
what explicit models can resolve.

\subsection{Lessons for FRFP and Long-Horizon Autonomy}

Waymo’s driverless service illustrates several key FRFP themes:

\begin{enumerate}
  \item \textbf{Explicit autonomy can work within strong structural
        constraints.}
        Carefully designed explicit semantics (perception, planning),
        validated in limited ODDs, can support safe autonomous operation
        for many trips and miles.

  \item \textbf{ODDs are practical safe horizons.}
        Geographic and situational constraints act as explicit safe horizons
        on where the autonomy stack is authorised to operate, aligning with
        FRFP’s call for bounded runs in high-stakes domains.

  \item \textbf{Incidents must feed into tacit governance.}
        Logging, simulation, and human safety reviews turn each anomaly
        into an opportunity to update both models and institutional
        understanding, rather than trusting explicit systems blindly.

  \item \textbf{Institutional non-automation remains.}
        Decisions about expansion to new cities, policies for interacting
        with emergency services, and responses to public concern cannot be
        automated; they require human and institutional judgement.

  \item \textbf{Formal FRFP perspective helps generalise autonomy design.}
        By seeing Waymo not just as a technical stack but as an FRFP system
        with explicit/tacit and institutional layers, we gain a template for
        evaluating and designing other long-horizon autonomous services
        (delivery robots, industrial fleets, etc.).
\end{enumerate}

In this sense, Waymo is a high-profile example of explicit AI autonomy that
has achieved non-trivial real-world deployment, while still depending
essentially on FRFP-style boundaries, safe horizons, and institutional
governance to remain acceptable and safe over long horizons.
